\section{Konu 2: Doğru ve Düzlem Denklemleri, Silindirler ve Kuadratik Yüzeyler}

Bu bölümde Konu~2 kapsamında derste (01.10, 02.10, 09.10, 15.10, 22.10, 23.10 ve 08.11 tarihli tahta/defter fotoğrafları) çözülen bütün örnekler, çözümsüz bırakılan çalışma soruları ve kaynak kitaplara benzer ek alıştırmalar adım adım çözülmüştür. Aynı örnek birden çok derste tekrarlandıysa bir kez yazılmış, tüm fotoğraf kaynakları verilmiştir. Sorular, konunun işleniş sırasına (doğru $\to$ düzlem $\to$ doğru--düzlem $\to$ yüzeyler) göre dizilmiştir.

\begin{bilgi}[Kullanılan temel formüller]
Doğru: $\frac{x-x_0}{a}=\frac{y-y_0}{b}=\frac{z-z_0}{c}$, $\mathbf v=(a,b,c)$.\quad
Düzlem: $Ax+By+Cz+D=0$, $\mathbf n=(A,B,C)$.\quad
Nokta--doğru: $\ell=\frac{\|\overrightarrow{AP}\times\mathbf v\|}{\|\mathbf v\|}$.\quad
Nokta--düzlem: $d=\frac{|Ax_0+By_0+Cz_0+D|}{\sqrt{A^2+B^2+C^2}}$.
\end{bilgi}

%======================================================================
\subsection{Uzayda doğru denklemleri}

\begin{soru}[Soru 1 \foto{IMG-20261009-WA0030}, \foto{IMG-20261023-WA0018}]
$A(1,2,3)$ noktasından geçen ve $\mathbf v=(2,3,0)$ vektörüne paralel olan doğrunun denklemini yazınız.
\end{soru}
\begin{cevap}[Çözüm (adım adım)]
\textbf{1. Verileri okuyalım.} Nokta $(x_0,y_0,z_0)=(1,2,3)$, doğrultman parametreleri $a=2$, $b=3$, $c=0$.

\textbf{2. Parametrik denklemler.} $x=1+2t,\quad y=2+3t,\quad z=3+0\cdot t=3$.

\textbf{3. Kartezyen biçim.} $c=0$ olduğundan $\frac{z-3}{0}$ yazılamaz; $z$ sabittir. İlk iki denklemden $t$ çekilip eşitlenir:
\[ \boxed{\frac{x-1}{2}=\frac{y-2}{3},\qquad z=3}. \]
\textbf{4. Sadeleştirme.} $3(x-1)=2(y-2)\Rightarrow3x-2y+1=0$. Doğru $\{3x-2y+1=0,\ z=3\}$ sistemiyle verilir.

\textit{Not:} Tahtada sonuç yalnız $3x-2y+1=0$ olarak kutulanmıştı; bu tek başına $Oz$'ye paralel bir düzlemdir. Doğru için $z=3$ de mutlaka yazılmalıdır.
\end{cevap}

\begin{soru}[Soru 2 \foto{IMG-20261009-WA0030}, \foto{IMG-20261023-WA0017}]
$A(1,3,-1)$ ve $B(3,5,6)$ noktalarından geçen doğrunun denklemini bulunuz.
\end{soru}
\begin{cevap}[Çözüm (adım adım)]
\textbf{1. Doğrultman.} $\overrightarrow{AB}=(3-1,\ 5-3,\ 6-(-1))=(2,2,7)$.

\textbf{2. Denklem.} $A$ noktası ve $\overrightarrow{AB}$ ile:
\[ \frac{x-1}{3-1}=\frac{y-3}{5-3}=\frac{z-(-1)}{6-(-1)}\ \Longrightarrow\ \boxed{\frac{x-1}{2}=\frac{y-3}{2}=\frac{z+1}{7}} \]
\textbf{3. Kontrol.} $B$'yi yerine koyalım: $\frac{3-1}{2}=1$, $\frac{5-3}{2}=1$, $\frac{6+1}{7}=1$; üç oran eşit \checkmark. Parametrik biçim $x=1+2t$, $y=3+2t$, $z=-1+7t$.
\end{cevap}

\begin{soru}[Soru 3 \foto{IMG-20261009-WA0030}, \foto{IMG-20261023-WA0017}]
$\dfrac{x-1}{3}=\dfrac{y+2}{m}=\dfrac{z-2}{4}$ ve $\dfrac{x-2}{6}=\dfrac{y-3}{8}=\dfrac{z+1}{n}$ doğruları paralel ise $m$ ve $n$ kaçtır?
\end{soru}
\begin{cevap}[Çözüm (adım adım)]
\textbf{1. Doğrultmanlar.} $\mathbf v=(3,m,4)$, $\mathbf u=(6,8,n)$.

\textbf{2. Paralellik şartı.} Doğrultmanlar orantılı olmalı: $\dfrac36=\dfrac m8=\dfrac4n$. İlk oran bilinen: $\frac36=\frac12$.

\textbf{3. Çözüm.} $\frac m8=\frac12\Rightarrow m=4$;\quad $\frac4n=\frac12\Rightarrow n=8$. Yani $\boxed{m=4,\ n=8}$.

\textbf{4. Çakışık mı?} $d_1$'in noktası $(1,-2,2)$: $\frac{1-2}{6}=-\frac16$, $\frac{-2-3}{8}=-\frac58$; oranlar eşit değil, nokta $d_2$'de değil. Doğrular paralel ve farklıdır.
\end{cevap}

%======================================================================
\subsection{İki doğru arasındaki açı ve noktanın doğruya uzaklığı}

\begin{soru}[Soru 4 \foto{IMG-20261002-WA0018}, \foto{IMG-20261023-WA0019}]
$\dfrac{x-1}{1}=\dfrac{y}{-\sqrt2}=\dfrac{z-1}{1}$ ve $\dfrac{x}{-1}=\dfrac{y-1}{\sqrt2}=\dfrac{z+2}{1}$ doğrularının oluşturduğu açıyı bulunuz.
\end{soru}
\begin{cevap}[Çözüm (adım adım)]
\textbf{1. Doğrultmanlar.} Paydalardan $\mathbf v=(1,-\sqrt2,1)$, $\mathbf u=(-1,\sqrt2,1)$.

\textbf{2. Skaler çarpım ve normlar.}
$\mathbf u\cdot\mathbf v=-1+(-\sqrt2)(\sqrt2)+1=-1-2+1=-2$;\quad $\|\mathbf u\|=\|\mathbf v\|=\sqrt{1+2+1}=2$.

\textbf{3. Açı.} $\cos\theta=\dfrac{-2}{2\cdot2}=-\dfrac12\Rightarrow\theta=120^\circ$ (doğrultman vektörleri arasındaki açı).

\textbf{4. Doğrular arasındaki açı.} Kesişen iki doğru $120^\circ$ ve $180^\circ-120^\circ=60^\circ$'lik açılar oluşturur; doğrular arasındaki açı dar açı olarak verilir: $\boxed{60^\circ}$. (Mutlak değerli formülle $\cos\theta=\frac{|-2|}{4}=\frac12$ doğrudan $60^\circ$ verir.)
\end{cevap}

\begin{soru}[Soru 5 \foto{IMG-20261002-WA0019}, \foto{IMG-20261023-WA0020}, \foto{IMG-20261023-WA0057}]
$P(2,3,-1)$ noktasının $\dfrac{x-1}{1}=\dfrac{y-2}{1}=\dfrac{z-13}{4}$ doğrusuna uzaklığını bulunuz.
\end{soru}
\begin{cevap}[Çözüm (adım adım)]
\textbf{1. Formül.} $A\in d$, $\mathbf v$ doğrultman ise $\ell=\dfrac{\|\overrightarrow{AP}\times\mathbf v\|}{\|\mathbf v\|}$ (paralelkenar alanı / taban).

\textbf{2. Nokta ve doğrultman.} $A(1,2,13)$, $\mathbf v=(1,1,4)$.

\textbf{3. $\overrightarrow{AP}$.} $(2-1,\ 3-2,\ -1-13)=(1,1,-14)$.

\textbf{4. Vektörel çarpım.}
\[ \overrightarrow{AP}\times\mathbf v=\begin{vmatrix}\mathbf i&\mathbf j&\mathbf k\\1&1&-14\\1&1&4\end{vmatrix}
=\mathbf i\,(1\cdot4-(-14)\cdot1)-\mathbf j\,(1\cdot4-(-14)\cdot1)+\mathbf k\,(1\cdot1-1\cdot1)=(18,-18,0). \]
\textbf{5. Uzaklık.}
\[ \ell=\frac{\sqrt{18^2+(-18)^2+0}}{\sqrt{1+1+16}}=\frac{18\sqrt2}{\sqrt{18}}=\frac{18\sqrt2}{3\sqrt2}=\boxed{6}. \]
\textit{Not:} Defterde önce $A(2,3,17)$ yazılmıştı; bu da doğrunun noktasıdır ($t=1$) ve aynı sonucu verir: $\overrightarrow{AP}=(0,0,-18)$, $\overrightarrow{AP}\times\mathbf v=(18,-18,0)$.
\end{cevap}

%======================================================================
\subsection{Düzlem denklemleri}

\begin{soru}[Soru 6 \foto{IMG-20261002-WA0020}, \foto{IMG-20261023-WA0030}]
$A(2,3,-5)$ noktasından geçen ve $\mathbf n=(7,3,4)$ vektörüne dik olan düzlemin denklemini bulunuz.
\end{soru}
\begin{cevap}[Çözüm (adım adım)]
\textbf{1.} Düzlem $\mathbf n$'ye dik $\Rightarrow$ $\mathbf n$ düzlemin normalidir: $A=7,\ B=3,\ C=4$.

\textbf{2.} $A(x-x_0)+B(y-y_0)+C(z-z_0)=0$:\quad $7(x-2)+3(y-3)+4(z+5)=0$.

\textbf{3.} $7x-14+3y-9+4z+20=0\Rightarrow\boxed{7x+3y+4z-3=0}$.

\textbf{4. Kontrol.} $14+9-20-3=0$ \checkmark.
\end{cevap}

\begin{soru}[Soru 7 \foto{IMG-20261002-WA0024}, \foto{IMG-20261002-WA0021}, \foto{IMG-20261023-WA0032}]
$P(2,3,5)$ noktasından geçen ve $\dfrac{x-1}{4}=\dfrac{y-2}{-1}=\dfrac{z+1}{6}$ doğrusuna dik olan düzlemin denklemini bulunuz.
\end{soru}
\begin{cevap}[Çözüm (adım adım)]
\textbf{1. Fikir.} Düzlem doğruya dik ise doğrunun doğrultmanı düzleme diktir, yani normale paraleldir: $\mathbf n=\mathbf v$.

\textbf{2.} $\mathbf v=(4,-1,6)=\mathbf n$, nokta $(2,3,5)$.

\textbf{3.} $4(x-2)-1(y-3)+6(z-5)=0\Rightarrow4x-8-y+3+6z-30=0\Rightarrow\boxed{4x-y+6z-35=0}$.

\textbf{4. Kontrol.} $8-3+30-35=0$ \checkmark.

\textit{Tahtadaki hata:} 23.10 tahtasında sonuç $4x-y+6z-3=0$ yazılmıştı. Sabit terim $-8+3-30=-35$'tir; $-3$ ile $P$ noktası denklemi sağlamaz ($8-3+30-3=32\neq0$). 02.10 tahtasında doğru sonuç ($-35$) yazılmıştı.
\end{cevap}

\begin{soru}[Soru 8 \foto{IMG-20261002-WA0022}, \foto{IMG-20261023-WA0050}]
$A(1,3,5)$, $B(4,2,1)$, $C(-1,2,-3)$ noktalarından geçen düzlemin denklemini bulunuz.
\end{soru}
\begin{cevap}[Çözüm (adım adım)]
\textbf{1. Fikir.} $P(x,y,z)$ düzlemde $\iff$ $\overrightarrow{AP},\overrightarrow{AB},\overrightarrow{AC}$ aynı düzlemde $\iff$ karma çarpım (determinant) $0$.

\textbf{2. Vektörler.} $\overrightarrow{AB}=(3,-1,-4)$, $\overrightarrow{AC}=(-2,-1,-8)$, $\overrightarrow{AP}=(x-1,y-3,z-5)$.

\textbf{3. Determinant.}
\[ \begin{vmatrix}x-1&y-3&z-5\\3&-1&-4\\-2&-1&-8\end{vmatrix}=0 . \]
Birinci satıra göre açılım (eş çarpanlar):
\begin{align*}
(x-1)\cdot\big[(-1)(-8)-(-4)(-1)\big]&=(x-1)(8-4)=4(x-1),\\
-(y-3)\cdot\big[3(-8)-(-4)(-2)\big]&=-(y-3)(-24-8)=32(y-3),\\
(z-5)\cdot\big[3(-1)-(-1)(-2)\big]&=(z-5)(-3-2)=-5(z-5).
\end{align*}
\textbf{4. Toplam.} $4(x-1)+32(y-3)-5(z-5)=0\Rightarrow4x+32y-5z-4-96+25=0\Rightarrow\boxed{4x+32y-5z-75=0}$.

\textbf{5. Sağlama.} $A$: $4+96-25-75=0$; $B$: $16+64-5-75=0$; $C$: $-4+64+15-75=0$ \checkmark.
\end{cevap}

%======================================================================
\subsection{İki düzlemin durumu ve noktanın düzleme uzaklığı}

\begin{soru}[Soru 9 \foto{IMG-20261002-WA0022}, \foto{IMG-20261023-WA0024}]
$2x+my+3z+4=0$ ve $-nx+4y+z+2=0$ düzlemleri birbirine dik ise $m$ ile $n$ arasındaki bağıntıyı bulunuz.
\end{soru}
\begin{cevap}[Çözüm (adım adım)]
\textbf{1.} Normaller $\mathbf n_1=(2,m,3)$, $\mathbf n_2=(-n,4,1)$.

\textbf{2.} Düzlemler dik $\iff$ normaller dik $\iff$ $A_1A_2+B_1B_2+C_1C_2=0$:
\[ 2(-n)+m\cdot4+3\cdot1=0\ \Longrightarrow\ \boxed{4m-2n+3=0}. \]
Bu bir tek çözüm değil, $m$ ile $n$ arasında bir bağıntıdır: $n=2m+\frac32$. Örneğin $m=1$ için $n=\frac72$: $(2,1,3)\cdot(-\frac72,4,1)=-7+4+3=0$ \checkmark.
\end{cevap}

\begin{soru}[Soru 10 \foto{IMG-20261002-WA0023}, \foto{IMG-20261023-WA0023}]
\textbf{(a)} $x+y+z=0$ ile $2x+2y+2z-5=0$ düzlemlerinin paralel mi, çakışık mı olduğunu belirleyiniz.
\textbf{(b)} $x+y+z=0$ ile $x-y=0$ düzlemlerinin dik olup olmadığını belirleyiniz.
\end{soru}
\begin{cevap}[Çözüm (adım adım)]
\textbf{(a)} $A_1=B_1=C_1=1$, $D_1=0$; $A_2=B_2=C_2=2$, $D_2=-5$.
\[ \frac{A_1}{A_2}=\frac{B_1}{B_2}=\frac{C_1}{C_2}=\frac12\ \Rightarrow\ \text{paralel};\qquad \frac{D_1}{D_2}=\frac{0}{-5}=0\neq\frac12\ \Rightarrow\ \text{çakışık değil.} \]
Düzlemler paralel ve farklıdır. (Aralarındaki uzaklık: ikinciyi $x+y+z-\frac52=0$ yazarsak $d=\frac{|{-\frac52}-0|}{\sqrt3}=\frac{5}{2\sqrt3}$.)

\textit{Not:} 02.10 tahtasında ikinci düzlem $2x+2y+2z=5$ biçiminde yazılıp kontrol $\frac12\neq\frac05$ diye yapılmıştı; $Ax+By+Cz+D=0$ biçiminde $D_2=-5$'tir. Sonuç değişmez.

\textbf{(b)} $\mathbf n_1=(1,1,1)$, $\mathbf n_2=(1,-1,0)$: $\mathbf n_1\cdot\mathbf n_2=1-1+0=0$ $\Rightarrow$ düzlemler \textbf{diktir}.
\end{cevap}

\begin{soru}[Soru 11 \foto{IMG-20261002-WA0023}, \foto{IMG-20261023-WA0025}, \foto{IMG-20261023-WA0031}]
$P(2,3,4)$ noktasının $3x-4y+12z+10=0$ düzlemine uzaklığını bulunuz.
\end{soru}
\begin{cevap}[Çözüm (adım adım)]
\textbf{1.} $A=3$, $B=-4$, $C=12$, $D=10$; $(x_0,y_0,z_0)=(2,3,4)$.

\textbf{2. Pay.} $|3\cdot2+(-4)\cdot3+12\cdot4+10|=|6-12+48+10|=52$.

\textbf{3. Payda.} $\sqrt{3^2+(-4)^2+12^2}=\sqrt{9+16+144}=\sqrt{169}=13$.

\textbf{4.} $d=\dfrac{52}{13}=\boxed{4}$.
\end{cevap}

\begin{soru}[Soru 12 \foto{IMG-20261002-WA0026}, \foto{IMG-20261023-WA0056}]
$P_1:\ Ax+By+Cz+D_1=0$ ve $P_2:\ Ax+By+Cz+D_2=0$ paralel düzlemleri arasındaki uzaklığın $d=\dfrac{|D_2-D_1|}{\sqrt{A^2+B^2+C^2}}$ olduğunu gösteriniz.
\end{soru}
\begin{cevap}[Çözüm (adım adım)]
\textbf{1.} Paralel düzlemlerde $P_1$'in her noktası $P_2$'ye aynı uzaklıktadır; $P_1$'den bir $K(x_0,y_0,z_0)$ alalım.

\textbf{2.} Noktanın düzleme uzaklığı: $d=\dfrac{|Ax_0+By_0+Cz_0+D_2|}{\sqrt{A^2+B^2+C^2}}$.

\textbf{3.} $K\in P_1$ olduğundan $Ax_0+By_0+Cz_0=-D_1$. Yerine yazarsak
\[ d=\frac{|-D_1+D_2|}{\sqrt{A^2+B^2+C^2}}=\frac{|D_2-D_1|}{\sqrt{A^2+B^2+C^2}} . \]
\textit{Uyarı:} Formül yalnızca iki denklem \emph{aynı} $A,B,C$ katsayılarıyla yazıldığında geçerlidir.
\end{cevap}

%======================================================================
\subsection{Doğru ile düzlem; arakesit doğrusu}

\begin{soru}[Soru 13 \foto{IMG-20261023-WA0023}, \foto{IMG-20261023-WA0025}]
$d:\ \dfrac{x-1}{2}=\dfrac{y}{3}=\dfrac{z+1}{-1}$ doğrusunun
\textbf{(a)} $4x+6y-2z+5=0$ düzlemine,
\textbf{(b)} $x-y-z+7=0$ düzlemine göre durumunu inceleyiniz.
\end{soru}
\begin{cevap}[Çözüm (adım adım)]
Doğrunun doğrultmanı $\mathbf v=(2,3,-1)$, üzerindeki bir nokta $(1,0,-1)$.

\textbf{(a)} $\mathbf n=(4,6,-2)$. Diklik testi: $\frac24=\frac36=\frac{-1}{-2}=\frac12$; oranlar eşit $\Rightarrow$ $\mathbf v\parallel\mathbf n$ $\Rightarrow$ \textbf{doğru düzleme diktir}.

\textbf{(b)} $\mathbf n=(1,-1,-1)$. Oranlar $\frac21\neq\frac3{-1}$, dik değil. Paralellik testi: $\mathbf n\cdot\mathbf v=2-3+1=0$ $\Rightarrow$ $\mathbf v\perp\mathbf n$: doğru düzleme paralel \emph{ya da} düzlemin içinde. Noktayı deneyelim: $1-0-(-1)+7=9\neq0$, nokta düzlemde değil. Sonuç: \textbf{doğru düzleme paraleldir}, düzlemin içinde değildir.

\textit{Not:} 23.10 tahtasında şart ``$aA+b.A+cC=0$'' yazılmıştı; doğrusu $aA+bB+cC=0$.
\end{cevap}

\begin{soru}[Soru 14 \foto{IMG-20261002-WA0026}, \foto{IMG-20261002-WA0027}, \foto{IMG-20261023-WA0061}]
$\dfrac{x-1}{2}=\dfrac{y-1}{-1}=\dfrac{z}{3}$ doğrusunun $2x+3y+z+11=0$ düzlemini kestiği noktayı bulunuz.
\end{soru}
\begin{cevap}[Çözüm (adım adım)]
\textbf{1. Parametrik.} Ortak oran $t$: $x=1+2t,\ y=1-t,\ z=3t$.

\textbf{2. Düzlemde yerine.} $2(1+2t)+3(1-t)+3t+11=0\Rightarrow2+4t+3-3t+3t+11=0\Rightarrow4t+16=0\Rightarrow t=-4$.

\textbf{3. Nokta.} $x=1-8=-7$, $y=1+4=5$, $z=-12$: $\boxed{(-7,5,-12)}$.

\textbf{4. Sağlama.} Düzlem: $-14+15-12+11=0$ \checkmark. Doğru: $\frac{-8}{2}=\frac{4}{-1}=\frac{-12}{3}=-4$ \checkmark.
\end{cevap}

\begin{soru}[Soru 15 \foto{IMG-20261002-WA0026}, \foto{IMG-20261023-WA0033}]
$x+y+z=3$ ve $x-y+z=1$ düzlemlerinin arakesit doğrusunu bulunuz.
\end{soru}
\begin{cevap}[Çözüm (adım adım)]
\textbf{1. Doğrultman.} Arakesit doğrusu iki düzlemde de yatar, iki normale de diktir:
\[ \mathbf v=\mathbf n_1\times\mathbf n_2=\begin{vmatrix}\mathbf i&\mathbf j&\mathbf k\\1&1&1\\1&-1&1\end{vmatrix}
=\mathbf i(1\cdot1-1\cdot(-1))-\mathbf j(1\cdot1-1\cdot1)+\mathbf k(1\cdot(-1)-1\cdot1)=(2,0,-2). \]
Sadeleştirerek $\mathbf v=(1,0,-1)$ alalım.

\textbf{2. Ortak nokta.} $z=0$: $x+y=3$, $x-y=1$ $\Rightarrow$ $x=2$, $y=1$; $A(2,1,0)$.

\textbf{3. Doğru.} $x=2+t,\ y=1,\ z=-t$, yani
\[ \boxed{\frac{x-2}{1}=\frac{z}{-1},\qquad y=1}. \]
\textbf{4. Sağlama.} $(2+t)+1+(-t)=3$ \checkmark; $(2+t)-1+(-t)=1$ \checkmark (her $t$ için).
\end{cevap}

%======================================================================
\subsection{Çalışma soruları}

\begin{soru}[Çalışma 1 (derste çalışma sorusu olarak verildi) \foto{IMG-20261001-WA0014}, \foto{IMG-20261002-WA0027}, \foto{IMG-20261023-WA0048}]
$P(1,0,2)$ noktasının $x+2y-2z+4=0$ düzlemine uzaklığını bulunuz.
\end{soru}
\begin{cevap}
$A=1,\ B=2,\ C=-2,\ D=4$:
\[ d=\frac{|1\cdot1+2\cdot0-2\cdot2+4|}{\sqrt{1+4+4}}=\frac{|1+0-4+4|}{3}=\boxed{\frac13}. \]
\end{cevap}

\begin{soru}[Çalışma 2 (derste çalışma sorusu olarak verildi) \foto{IMG-20261001-WA0014}, \foto{IMG-20261002-WA0027}, \foto{IMG-20261023-WA0048}]
$A(1,0,0)$, $B(0,1,0)$, $C(0,0,1)$ noktalarından geçen düzlemin denklemini bulunuz.
\end{soru}
\begin{cevap}[Çözüm (adım adım)]
$\overrightarrow{AB}=(-1,1,0)$, $\overrightarrow{AC}=(-1,0,1)$, $\overrightarrow{AP}=(x-1,y,z)$.
\[ \begin{vmatrix}x-1&y&z\\-1&1&0\\-1&0&1\end{vmatrix}=(x-1)(1-0)-y(-1-0)+z(0+1)=x-1+y+z=0 . \]
$\boxed{x+y+z=1}$. Kontrol: üç nokta da sağlar \checkmark. Başka yol: eksenleri $1,1,1$'de kesen düzlem, eksen kesişim biçiminden $\frac x1+\frac y1+\frac z1=1$.
\end{cevap}

\begin{soru}[Çalışma 3 (derste çalışma sorusu olarak verildi) \foto{IMG-20261001-WA0014}, \foto{IMG-20261002-WA0027}, \foto{IMG-20261023-WA0048}]
$\dfrac{x-2}{1}=\dfrac{y}{2}=\dfrac{z+1}{-1}$ doğrusu ile $x-y+2z=3$ düzleminin kesim noktasını bulunuz.
\end{soru}
\begin{cevap}[Çözüm (adım adım)]
Parametrik: $x=2+t,\ y=2t,\ z=-1-t$. Düzlemde yerine:
\[ (2+t)-2t+2(-1-t)=3\ \Rightarrow\ 2+t-2t-2-2t=3\ \Rightarrow\ -3t=3\ \Rightarrow\ t=-1 . \]
Nokta: $x=1$, $y=-2$, $z=0$: $\boxed{(1,-2,0)}$. Sağlama: $1-(-2)+0=3$ \checkmark.
\end{cevap}

\begin{soru}[Çalışma 4 (derste çalışma sorusu olarak verildi) \foto{IMG-20261001-WA0014}, \foto{IMG-20261002-WA0027}, \foto{IMG-20261023-WA0048}]
$\dfrac{x}{1}=\dfrac{y}{1}=\dfrac{z}{0}$ ve $\dfrac{x-1}{1}=\dfrac{y}{-1}=\dfrac{z}{0}$ doğruları arasındaki açıyı bulunuz.
\end{soru}
\begin{cevap}[Çözüm (adım adım)]
\textbf{1. Yazımın anlamı.} Paydadaki $0$ bir bölme değildir; doğrultmanın üçüncü bileşeni $0$ demektir, yani $z=0$. Doğrular: $d_1:\ x=t,\ y=t,\ z=0$ ve $d_2:\ x=1+s,\ y=-s,\ z=0$.

\textbf{2. Doğrultmanlar.} $\mathbf v=(1,1,0)$, $\mathbf u=(1,-1,0)$; $\mathbf u\cdot\mathbf v=1-1+0=0$.

\textbf{3. Sonuç.} $\cos\theta=0\Rightarrow\boxed{\theta=90^\circ}$: doğrular diktir. (Üstelik ikisi de $z=0$ düzleminde olup $t=s+1$, $t=-s$ $\Rightarrow$ $s=-\frac12$, $t=\frac12$ için $(\frac12,\frac12,0)$ noktasında kesişirler.)
\end{cevap}

\begin{soru}[Çalışma 5 (derste çalışma sorusu olarak verildi) \foto{IMG-20261009-WA0032}, \foto{IMG-20261023-WA0027}]
$(2,-1,3)$ ile $(1,4,2)$ vektörlerinin dik olup olmadığını belirleyiniz ve aralarındaki açının kosinüsünü bulunuz.
\end{soru}
\begin{cevap}
(Bu soru Konu~1 bilgisiyle çözülür; derste Konu~2 çalışma sorularıyla birlikte verildi.)
$(2,-1,3)\cdot(1,4,2)=2-4+6=4\neq0$ $\Rightarrow$ \textbf{dik değildir}. $\|(2,-1,3)\|=\sqrt{14}$, $\|(1,4,2)\|=\sqrt{21}$:
\[ \cos\theta=\frac{4}{\sqrt{14}\sqrt{21}}=\frac{4}{\sqrt{294}}=\frac{4}{7\sqrt6}=\frac{2\sqrt6}{21}\approx0{,}233\qquad(\theta\approx76{,}5^\circ). \]
\end{cevap}

\begin{soru}[Çalışma 6 (derste çalışma sorusu olarak verildi) \foto{IMG-20261009-WA0032}, \foto{IMG-20261023-WA0027}]
$x^2+y^2+z^2+4x-6y+2z-11=0$ küresinin merkezini ve yarıçapını bulunuz.
\end{soru}
\begin{cevap}
Kareye tamamlama: $(x^2+4x+4)+(y^2-6y+9)+(z^2+2z+1)=11+4+9+1$, yani
\[ (x+2)^2+(y-3)^2+(z+1)^2=25 . \]
Merkez $\boxed{M(-2,3,-1)}$, yarıçap $\boxed{r=5}$. (Küre, elipsoidin $a=b=c$ özel hâlidir.)
\end{cevap}

\begin{soru}[Çalışma 7 (derste çalışma sorusu olarak verildi) \foto{IMG-20261009-WA0032}, \foto{IMG-20261023-WA0027}]
$A(0,1,2)$ ve $B(3,1,-1)$ noktalarından geçen doğrunun parametrik ve kartezyen denklemlerini bulunuz.
\end{soru}
\begin{cevap}[Çözüm (adım adım)]
\textbf{1.} $\overrightarrow{AB}=(3,0,-3)\parallel(1,0,-1)$.

\textbf{2. Parametrik.} $A$ ve $\overrightarrow{AB}$ ile: $x=3t,\ y=1,\ z=2-3t$ ($t\in\mathbb R$). (Sadeleşmiş doğrultmanla: $x=t,\ y=1,\ z=2-t$.)

\textbf{3. Kartezyen.} $y$ bileşeni $0$ olduğundan $y=1$ sabit:
\[ \boxed{\frac{x}{3}=\frac{z-2}{-3},\qquad y=1}\quad(\text{yani } x+z=2,\ y=1). \]
Kontrol: $B$: $\frac33=1$, $\frac{-1-2}{-3}=1$, $y=1$ \checkmark.
\end{cevap}

%======================================================================
\subsection{Silindirler ve kuadratik yüzeyler}

\begin{soru}[Soru 16 \foto{IMG-20261022-WA0050}, \foto{IMG-20261022-WA0052}, \foto{IMG-20261108-WA0096}]
Tahtada örnek olarak verilen $x^2+y^2=a^2$, $z=4x^2$ ve $-\dfrac{x^2}{a^2}+\dfrac{y^2}{b^2}=1$ yüzeylerinin neden silindir olduğunu açıklayınız ve her birinin ana doğrularının yönünü ve dayanak eğrisini belirtiniz.
\end{soru}
\begin{cevap}[Çözüm (adım adım)]
Bir denklemde değişkenlerden biri yoksa, o değişken serbesttir: dayanak eğrisi üzerindeki her noktadan o eksene paralel doğru tamamen yüzeyde kalır.
\begin{itemize}
  \item $x^2+y^2=a^2$: $z$ yok $\Rightarrow$ ana doğrular $Oz$'ye paralel; dayanak eğrisi $xy$-düzleminde yarıçapı $a$ olan çember $\Rightarrow$ \textbf{dairesel silindir}.
  \item $z=4x^2$: $y$ yok $\Rightarrow$ ana doğrular $Oy$'ye paralel; dayanak eğrisi $xz$-düzleminde $z=4x^2$ parabolü $\Rightarrow$ \textbf{parabolik silindir} ($y$ ekseni boyunca uzanan oluk).
  \item $-\frac{x^2}{a^2}+\frac{y^2}{b^2}=1$: $z$ yok $\Rightarrow$ ana doğrular $Oz$'ye paralel; dayanak eğrisi $xy$-düzleminde hiperbol. $\frac{y^2}{b^2}=1+\frac{x^2}{a^2}\ge1$ olduğundan $|y|\ge b$: hiperbolün kolları $\pm y$ yönünde açılır ($y=\pm b$ tepe noktaları). Yüzey $y$ ekseninin iki yanında, $z$ yönünde uzanan iki ayrı yapraktan oluşan \textbf{hiperbolik silindir}dir.
\end{itemize}
\end{cevap}

\begin{soru}[Soru 17 \foto{IMG-20261022-WA0049}, \foto{IMG-20261108-WA0097}]
Tek kanatlı hiperboloid $\frac{x^2}{a^2}+\frac{y^2}{b^2}-\frac{z^2}{c^2}=1$ için ``$z=0$'da elips, $x=0$'da hiperbol'' olduğunu; çift kanatlı hiperboloid $\frac{x^2}{a^2}+\frac{y^2}{b^2}-\frac{z^2}{c^2}=-1$ için de ``$-c<z<c$ aralığında yüzey yok'' olduğunu gösteriniz.
\end{soru}
\begin{cevap}[Çözüm (adım adım)]
\textbf{Tek kanatlı.} $z=0$ koyunca $\frac{x^2}{a^2}+\frac{y^2}{b^2}=1$: yarı eksenleri $a,b$ olan \textbf{elips} (bel elipsi). $x=0$ koyunca $\frac{y^2}{b^2}-\frac{z^2}{c^2}=1$: $yz$-düzleminde \textbf{hiperbol}. Genel $z=k$ izi $\frac{x^2}{a^2}+\frac{y^2}{b^2}=1+\frac{k^2}{c^2}>0$, her $k$ için elips olduğundan yüzey tek parçadır.

\textbf{Çift kanatlı.} $z=k$ koyunca
\[ \frac{x^2}{a^2}+\frac{y^2}{b^2}=\frac{k^2}{c^2}-1 . \]
Sol taraf $\ge0$ olduğundan $\frac{k^2}{c^2}\ge1$, yani $|k|\ge c$ olmalıdır. $-c<k<c$ ise sağ taraf negatiftir ve hiçbir $(x,y)$ denklemi sağlamaz: bu şeritte yüzeyin noktası yoktur. $|k|=c$'de tek nokta $(0,0,\pm c)$ (tepe noktaları), $|k|>c$'de elipsler vardır. Böylece yüzey $z\ge c$ ve $z\le-c$'de iki ayrı kanattan oluşur.
\end{cevap}

\begin{soru}[Soru 18 \foto{IMG-20261015-WA0043}, \foto{IMG-20261023-WA0063}]
Elipsoid $\frac{x^2}{a^2}+\frac{y^2}{b^2}+\frac{z^2}{c^2}=1$'in eksenleri $(\mp a,0,0)$, $(0,\mp b,0)$, $(0,0,\mp c)$ noktalarında kestiğini gösteriniz ve üst/alt yarı elipsoidin denklemlerini yazınız.
\end{soru}
\begin{cevap}
$Ox$ ekseninde $y=z=0$: $\frac{x^2}{a^2}=1\Rightarrow x=\mp a$. Benzer şekilde $Oy$ üzerinde $y=\mp b$, $Oz$ üzerinde $z=\mp c$.
$z$ çekilirse $z^2=c^2\left(1-\frac{x^2}{a^2}-\frac{y^2}{b^2}\right)$, yani
\[ z=c\sqrt{1-\frac{x^2}{a^2}-\frac{y^2}{b^2}}\ \ (\text{üst yarı}),\qquad z=-c\sqrt{1-\frac{x^2}{a^2}-\frac{y^2}{b^2}}\ \ (\text{alt yarı}), \]
tanım kümesi $\frac{x^2}{a^2}+\frac{y^2}{b^2}\le1$ elips bölgesidir.
\end{cevap}

\begin{soru}[Soru 19 \foto{IMG-20261022-WA0052}, \foto{IMG-20261108-WA0096}, \foto{IMG-20261108-WA0095}]
\textbf{(a)} $\frac{z^2}{c^2}=\frac{x^2}{a^2}+\frac{y^2}{b^2}$ konisinin $x=0$ izinin bir doğru çifti olduğunu gösteriniz.
\textbf{(b)} $\frac zc=\frac{x^2}{a^2}-\frac{y^2}{b^2}$ ($c>0$) eyer yüzeyinin $y=0$ ve $x=0$ izlerini bulunuz.
\end{soru}
\begin{cevap}
\textbf{(a)} $x=0$: $\frac{z^2}{c^2}=\frac{y^2}{b^2}\Rightarrow z=\pm\frac cb\,y$; orijinden geçen iki doğru. Koni bu tür doğrulardan (tepeden geçen ana doğrulardan) oluşur.

\textbf{(b)} $y=0$: $z=\frac{c}{a^2}x^2$, yukarı açılan parabol. $x=0$: $z=-\frac{c}{b^2}y^2$, aşağı açılan parabol. Orijin, bir yönde en alçak, diğer yönde en yüksek noktadır: \textbf{eyer noktası}. $z=0$ izi $y=\pm\frac ba x$ doğru çiftidir.
\end{cevap}

%======================================================================
\subsection{Ek alıştırmalar}

\begin{soru}[\ekal Ek alıştırma 1]
$d_1:\ \mathbf r=(1,1,0)+t(1,0,1)$ ve $d_2:\ \mathbf r=(0,2,1)+s(0,1,1)$ doğrularının durumunu belirleyiniz; aykırı iseler aralarındaki uzaklığı bulunuz.
\end{soru}
\begin{cevap}[Çözüm (adım adım)]
\textbf{1. Paralel mi?} $(1,0,1)$ ile $(0,1,1)$ orantılı değil $\Rightarrow$ paralel değil.

\textbf{2. Kesişiyor mu?} $x$: $1+t=0\Rightarrow t=-1$; $y$: $1=2+s\Rightarrow s=-1$; $z$: $t=-1$ fakat $1+s=0$; $-1\neq0$ $\Rightarrow$ kesişmez. Doğrular \textbf{aykırıdır}.

\textbf{3. Ortak dik.} $\mathbf v\times\mathbf u=\begin{vmatrix}\mathbf i&\mathbf j&\mathbf k\\1&0&1\\0&1&1\end{vmatrix}=(0-1,\ -(1-0),\ 1-0)=(-1,-1,1)$.

\textbf{4. Uzaklık.} $\overrightarrow{P_1P_2}=(0,2,1)-(1,1,0)=(-1,1,1)$;
\[ d=\frac{|(-1)(-1)+(1)(-1)+(1)(1)|}{\sqrt{1+1+1}}=\frac{1}{\sqrt3}=\frac{\sqrt3}{3}. \]
\end{cevap}

\begin{soru}[\ekal Ek alıştırma 2]
$\dfrac{x-1}{2}=\dfrac{y+2}{1}=\dfrac{z}{3}$ doğrusunu ve $Q(0,1,1)$ noktasını içeren düzlemin denklemini bulunuz.
\end{soru}
\begin{cevap}[Çözüm (adım adım)]
\textbf{1.} Doğrunun noktası $A(1,-2,0)$, doğrultmanı $\mathbf v=(2,1,3)$; $\overrightarrow{AQ}=(-1,3,1)$. Her iki vektör de düzlemde yatar.

\textbf{2. Normal.}
\[ \mathbf n=\mathbf v\times\overrightarrow{AQ}=\begin{vmatrix}\mathbf i&\mathbf j&\mathbf k\\2&1&3\\-1&3&1\end{vmatrix}=(1-9,\ -(2+3),\ 6+1)=(-8,-5,7). \]
\textbf{3. Düzlem.} $-8(x-1)-5(y+2)+7(z-0)=0\Rightarrow-8x-5y+7z-2=0\Rightarrow\boxed{8x+5y-7z+2=0}$.

\textbf{4. Kontrol.} $A$: $8-10-0+2=0$; $Q$: $0+5-7+2=0$; doğrunun başka bir noktası $(3,-1,3)$: $24-5-21+2=0$ \checkmark.
\end{cevap}

\begin{soru}[\ekal Ek alıştırma 3]
$2x-y+2z=3$ ile $-4x+2y-4z=6$ düzlemlerinin paralel olduğunu gösterip aralarındaki uzaklığı bulunuz.
\end{soru}
\begin{cevap}[Çözüm (adım adım)]
\textbf{1.} $\frac{2}{-4}=\frac{-1}{2}=\frac{2}{-4}=-\frac12$ $\Rightarrow$ paralel. İkinci denklemi $-2$'ye bölersek $2x-y+2z=-3$; birincisiyle aynı değil $\Rightarrow$ çakışık değil.

\textbf{2.} Aynı katsayılarla: $P_1:\ 2x-y+2z-3=0$ ($D_1=-3$), $P_2:\ 2x-y+2z+3=0$ ($D_2=3$).

\textbf{3.} $d=\dfrac{|3-(-3)|}{\sqrt{4+1+4}}=\dfrac63=\boxed{2}$.
\end{cevap}

\begin{soru}[\ekal Ek alıştırma 4]
$2x-y+z=4$ ve $x+y+2z=1$ düzlemleri arasındaki açıyı ve arakesit doğrusunu bulunuz.
\end{soru}
\begin{cevap}[Çözüm (adım adım)]
\textbf{1. Açı.} $\mathbf n_1=(2,-1,1)$, $\mathbf n_2=(1,1,2)$, $\mathbf n_1\cdot\mathbf n_2=2-1+2=3$, $\|\mathbf n_1\|=\|\mathbf n_2\|=\sqrt6$:
$\cos\theta=\frac{3}{6}=\frac12\Rightarrow\theta=60^\circ$.

\textbf{2. Doğrultman.}
$\mathbf n_1\times\mathbf n_2=\begin{vmatrix}\mathbf i&\mathbf j&\mathbf k\\2&-1&1\\1&1&2\end{vmatrix}=(-2-1,\ -(4-1),\ 2+1)=(-3,-3,3)\parallel(1,1,-1)$.

\textbf{3. Nokta.} $z=0$: $2x-y=4$, $x+y=1$ $\Rightarrow$ $3x=5$, $x=\frac53$, $y=-\frac23$.

\textbf{4. Doğru.} $x=\frac53+t,\ y=-\frac23+t,\ z=-t$. Kontrol: $2(\frac53+t)-(-\frac23+t)-t=\frac{12}{3}=4$ \checkmark; $(\frac53+t)+(-\frac23+t)-2t=1$ \checkmark.
\end{cevap}

\begin{soru}[\ekal Ek alıştırma 5]
$P(1,2,3)$ noktasının $x+y+z=0$ düzlemi üzerindeki dik izdüşümünü (dikme ayağını) ve noktanın düzleme uzaklığını bulunuz.
\end{soru}
\begin{cevap}[Çözüm (adım adım)]
\textbf{1.} $P$'den geçen ve normale $\mathbf n=(1,1,1)$ paralel doğru: $x=1+t,\ y=2+t,\ z=3+t$.

\textbf{2.} Düzlemde yerine: $(1+t)+(2+t)+(3+t)=0\Rightarrow6+3t=0\Rightarrow t=-2$.

\textbf{3.} Dikme ayağı $K(-1,0,1)$.

\textbf{4.} Uzaklık: $|PK|=\sqrt{4+4+4}=2\sqrt3$. Formülle: $\frac{|1+2+3|}{\sqrt3}=\frac{6}{\sqrt3}=2\sqrt3$ \checkmark.
\end{cevap}

\begin{soru}[\ekal Ek alıştırma 6]
$x^2+4y^2+z^2-2x+8y+1=0$ yüzeyini tanıyınız.
\end{soru}
\begin{cevap}[Çözüm (adım adım)]
\textbf{1. Kareye tamamlama.} $x^2-2x=(x-1)^2-1$, $4y^2+8y=4(y+1)^2-4$.

\textbf{2.} $(x-1)^2-1+4(y+1)^2-4+z^2+1=0\Rightarrow(x-1)^2+4(y+1)^2+z^2=4$.

\textbf{3.} $4$'e bölersek
\[ \frac{(x-1)^2}{4}+\frac{(y+1)^2}{1}+\frac{z^2}{4}=1 . \]
Üç kare, hepsi pozitif $\Rightarrow$ \textbf{elipsoid}; merkezi $(1,-1,0)$, yarı eksenleri $x$ yönünde $2$, $y$ yönünde $1$, $z$ yönünde $2$. ($a=c$ olduğundan $Oy$'ye paralel eksen etrafında dönel bir elipsoiddir.)
\end{cevap}

\begin{soru}[\ekal Ek alıştırma 7]
$y=z^2-x^2$ yüzeyini izleriyle tanıyınız.
\end{soru}
\begin{cevap}[Çözüm (adım adım)]
İki değişken kare, zıt işaretli; $y$ birinci dereceden $\Rightarrow$ \textbf{hiperbolik paraboloid} (eyer), ekseni $Oy$.
\begin{itemize}
  \item $y=k$: $z^2-x^2=k$; $k\neq0$ için hiperbol ($k>0$'da $z$ yönünde, $k<0$'da $x$ yönünde açılır), $k=0$ için $z=\pm x$ doğru çifti.
  \item $x=k$: $y=z^2-k^2$, $+y$ yönünde açılan parabol.
  \item $z=k$: $y=k^2-x^2$, $-y$ yönünde açılan parabol.
\end{itemize}
Orijin eyer noktasıdır: $z$ yönünde yürüyünce $y$ artar, $x$ yönünde yürüyünce azalır.
\end{cevap}
